% ============================================================
% PRN232 - Section 9: Conclusion (Vietnamese & English)
% ============================================================
\section{\bi{Kết luận và hướng phát triển}{Conclusion and Future Work}}
\label{sec:concl}

\subsection{\bi{Kết luận}{Conclusion}}

\bi{
Bài báo đã trình bày OmniCast --- một hệ thống lịch phát, truyền phát trực tiếp (LiveStream) và phân phối
nội dung số đa nền tảng được thực hiện bởi nhóm 4 sinh viên (\textbf{Nguyễn Bình An}, \textbf{Trần Thanh Sơn}, \textbf{Nguyễn Vĩnh Phú}, \textbf{Đồng Hoàng Nguyên}) tại \textbf{FPTU Quy Nhon AI Campus --- FU Cahata} dưới sự hướng dẫn của giảng viên \textbf{Thầy Huỳnh Vang}. Bốn đóng góp công
nghệ chính được chứng minh bằng mã nguồn và kiểm thử thực tế: (1) kiến trúc LiveStream \& Live Studio hai chiều kết hợp WebRTC WHIP/WHEP, MediaMTX HLS và Socket.IO \texttt{/live}; (2) thuật
toán phát hiện xung đột lịch phát bằng một truy vấn SQL thô tính
thời điểm kết thúc hiệu lực; (3) cơ chế điền lịch EPG 24/7 với lựa chọn
ngẫu nhiên có hạt giống bảo đảm tính tái lập; và (4) quy trình biên tập
AI gọi mô hình một lần duy nhất, kiểm tra schema bằng Zod, có đường
dự phòng heuristic và dấu vết phiên bản đầy đủ.

Toàn bộ hệ thống --- 19 module NestJS, hơn 140 endpoint REST, 21 model
Prisma, cổng web Next.js 15 với 33 route, ứng dụng di động Flutter 31 màn hình và microservice ghi nhật ký
kiểm toán gRPC --- đã được triển khai trên Vercel, kiểm thử bằng 482 kiểm
thử tự động trong quy trình tích hợp liên tục (đạt 100\%), và vận hành ổn định trên môi trường sản xuất.
}{
This paper has presented OmniCast, a multi-platform digital broadcast
scheduling, real-time LiveStream and delivery system developed by the four-member engineering team (\textbf{Nguyễn Bình An}, \textbf{Trần Thanh Sơn}, \textbf{Nguyễn Vĩnh Phú}, \textbf{Đồng Hoàng Nguyên}) at \textbf{FPTU Quy Nhon AI Campus --- FU Cahata} under the mentorship of \textbf{Mr. Huỳnh Vang}. Its four principal
technical contributions are demonstrated through the source code and
real test runs: (1) a bi-directional LiveStream and Live Studio architecture combining WebRTC WHIP/WHEP, MediaMTX HLS and Socket.IO \texttt{/live}; (2) a broadcast-conflict algorithm implemented as a single
raw SQL query that computes the effective end time; (3) a 24/7 EPG filler
whose seeded selection guarantees reproducibility; and (4) an AI
curation pipeline that issues a single consolidated model call,
validates the output against a Zod schema, offers a heuristic fallback
and records full version provenance.

The complete system --- 19 NestJS modules, 140+ REST endpoints, 21
Prisma models, a Next.js 15 portal with 33 routes, a 31-screen Flutter mobile app and a gRPC audit
microservice --- has been deployed on Vercel, validated by 482
automated tests (100\% pass rate) inside a continuous-integration pipeline, and operates
reliably in production.
}

\subsection{\bi{Hạn chế đã biết}{Known Limitations}}

\bi{
Ba hạn chế đã biết cần được nêu rõ. Thứ nhất, mức độ phủ kiểm thử
E2E trình duyệt cần mở rộng thêm cho các kịch bản đẩy luồng phần cứng WHIP cần máy chủ MediaMTX chuyên dụng. Thứ hai, lớp
ingest hiện hỗ trợ YouTube, Twitch và RSS theo chu kỳ cron. Thứ ba, thuật toán điền lịch chọn theo thể
loại và theo hạt giống, chưa tối ưu cá nhân hoá theo hành vi người xem thực tế.
}{
Three known limitations should be stated. First, browser E2E test coverage can be expanded further for hardware WHIP publishing scenarios requiring a dedicated MediaMTX instance in CI.
Second, the external ingestion layer currently polls YouTube, Twitch and RSS on cron intervals. Third, the filler algorithm selects by
category and by seed rather than by individualised viewer history.
}

\subsection{\bi{Hướng phát triển}{Future Work}}

\bi{
Bốn hướng phát triển được ưu tiên cho các phiên bản tiếp theo:
}{
Four directions are prioritised for subsequent iterations:
}

\begin{enumerate}[leftmargin=*,itemsep=1pt]
  \item \bi{\textbf{Cá nhân hoá lịch phát bằng học máy.} Thay vì
  chọn theo thể loại, thuật toán điền lịch có thể dùng lịch sử xem
  và dữ liệu tương tác để đề xuất nội dung cho từng kênh, đồng thời
  giữ nguyên tính chất xác định để có thể kiểm thử.}{
  \textbf{Recommender-driven scheduling.} Instead of selecting by
  category, the filler algorithm could use watch history and
  interaction data to propose content per channel while preserving
  determinism so the result remains testable.}

  \item \bi{\textbf{Mở rộng kiểm thử E2E MediaMTX.} Xây dựng môi trường
  MediaMTX trong CI để kiểm thử tự động luồng WHIP phát và HLS phát
  lại, đồng thời bổ sung kiểm thử trình duyệt cho vai trò STAFF và
  ADMIN.}{
  \textbf{Scaling MediaMTX end-to-end testing.} Provision a MediaMTX instance
  in CI to test the WHIP publish and HLS playback paths
  automatically, and add browser scenarios for the STAFF and ADMIN
  roles.}

  \item \bi{\textbf{Cải thiện trợ lý AI.} Bổ sung đánh giá nhiều mô
  hình, kiểm duyệt đầu ra bằng mô hình thứ hai và xếp hạng các báo
  cáo AI theo phản hồi của biên tập viên để tinh chỉnh prompt theo
  vòng lặp.}{
  \textbf{Improving the AI assistant.} Add multi-model evaluation,
  second-model output verification, and ranking of AI reports by
  editor feedback so the prompt can be tuned iteratively.}

  \item \bi{\textbf{Tách lớp truy cập dữ liệu.} Đưa truy vấn EPG vào
  một kho phân tích chỉ đọc với các chỉ mục bao phủ để tách tải
  đọc khỏi tải ghi trên cùng một cơ sở dữ liệu.}{
  \textbf{Separating the data-access tier.} Move EPG queries into a
  read-only analytical store with covering indexes to separate read
  load from write load on the same database.}
\end{enumerate}

\subsection{\bi{Lời cảm ơn và Tính tái lập}{Acknowledgements and Reproducibility Statement}}

\bi{
Nhóm tác giả (\textbf{Nguyễn Bình An}, \textbf{Trần Thanh Sơn}, \textbf{Nguyễn Vĩnh Phú}, \textbf{Đồng Hoàng Nguyên}) xin bày tỏ lòng biết ơn sâu sắc tới giảng viên hướng dẫn \textbf{Thầy Huỳnh Vang} cùng tập thể giảng viên tại \textbf{FPTU Quy Nhon AI Campus --- FU Cahata} đã tận tình định hướng kiến trúc, góp ý chuyên môn và đồng hành cùng nhóm trong suốt học kỳ Fall 2026. Toàn bộ mã nguồn và kịch bản kiểm thử được công bố tại kho lưu trữ \texttt{TomOutfit/OmniCast\_PRN232-PRM393\_Project\_Fa26Semester}.
}{
The authors (\textbf{Nguyễn Bình An}, \textbf{Trần Thanh Sơn}, \textbf{Nguyễn Vĩnh Phú}, \textbf{Đồng Hoàng Nguyên}) express their sincere gratitude to Mentor \textbf{Mr. Huỳnh Vang} and the faculty at \textbf{FPTU Quy Nhon AI Campus --- FU Cahata} for their architectural guidance and mentorship throughout the Fall 2026 semester. All source code and automated test suites are publicly reproducible in \texttt{TomOutfit/OmniCast\_PRN232-PRM393\_Project\_Fa26Semester}.
}
